
\documentclass{article}
\usepackage{colortbl}
\usepackage{makecell}
\usepackage{multirow}
\usepackage{supertabular}

\begin{document}

\newcounter{utterance}

\centering \large Interaction Transcript for game `clean\_up', experiment `2\_hard\_3obj\_ga', episode 3 with gemma{-}4{-}26b{-}a4b{-}it.
\vspace{24pt}

{ \footnotesize  \setcounter{utterance}{1}
\setlength{\tabcolsep}{0pt}
\begin{supertabular}{c@{$\;$}|p{.15\linewidth}@{}p{.15\linewidth}p{.15\linewidth}p{.15\linewidth}p{.15\linewidth}p{.15\linewidth}}
    \# & $\;$A & \multicolumn{4}{c}{Game Master} & $\;\:$B\\
    \hline

    \theutterance \stepcounter{utterance}  
    & & \multicolumn{4}{p{0.6\linewidth}}{
        \cellcolor[rgb]{0.9,0.9,0.9}{
            \makecell[{{p{\linewidth}}}]{
                \texttt{\tiny{[P1$\langle$GM]}}
                \texttt{Is mise do mháistir cluiche, agus tá tú ag imirt cluiche comhoibríoch leis an ngreille seo a leanas mar chlár cluiche:} \\
\\ 
\texttt{```} \\
\texttt{ 1234567} \\
\texttt{╔═╤═╤═══╗} \\
\texttt{║L│◌│◌◌◌║ 1} \\
\texttt{║C├─┤◌◌◌║ 2} \\
\texttt{╟─┤◌├─┬─╢ 3} \\
\texttt{║◌│◌│◌│◌║ 4} \\
\texttt{╟─┼─┘◌│◌║ 5} \\
\texttt{║◌│◌◌┌┴─╢ 6} \\
\texttt{║◌│◌◌│◌P║ 7} \\
\texttt{╚═╧══╧══╝} \\
\\ 
\texttt{```} \\
\\ 
\texttt{* Taispeánann an imeall uachtarach comhordanáidí x ag méadú i dtreo na láimhe deise, agus an t{-}imeall ar dheis comhordanáidí y ag méadú síos.} \\
\texttt{* Cuirtear na rudaí seo a leanas go randamach ar do ghreille: 'C', 'L', 'P'.} \\
\\ 
\\ 
\texttt{Feiceann an t{-}imreoir eile leagan éagsúil den chlár cluiche, áit a gcuirtear na rudaí in áiteanna randamacha éagsúla. Ní féidir leat clár an imreora eile a fheiceáil, agus ní féidir leo do cheannsa a fheiceáil.} \\
\\ 
\\ 
\texttt{**Sprioc:**} \\
\\ 
\texttt{Caithfidh an dá imreoir na rudaí a bhogadh ar a gcúlra féin ionas go gcríochnóidh rudaí comhionanna ag na comhordanáidí céanna. Caithfidh tú cumarsáid a dhéanamh leis an imreoir eile chun cumraíocht sprice choiteann a aontú.} \\
\\ 
\\ 
\texttt{**Rialacha:**} \\
\\ 
\texttt{* I ngach babhta, is féidir leat ceann amháin go díreach den dá ordú seo a leanas a sheoladh:} \\
\texttt{ 1. `ABAIR: <TEACHTAIREACHT>`: chun teachtaireacht (gach rud suas go dtí an chéad bhriseadh líne eile) a sheoladh chuig an imreoir eile. Cuirfidh mé ar aghaidh í chuig do pháirtí.} \\
\texttt{2. `BOG: <RUD>, (<X>, <Y>)`: chun rud a bhogadh go suíomh nua, áit a bhfuil `<X>` mar an colún agus `<Y>` mar an ró. Cuirfidh mé in iúl duit an raibh do bhogadh rathúil nó nach raibh.} \\
\texttt{* Mura gcloíonn tú leis an bhformáid, nó má sheolann tú roinnt orduithe ag an am céanna, caithfidh mé pionós a ghearradh ort.} \\
\texttt{* Má charnaíonn an dá imreoir níos mó ná 8 pionóis, caillfidh sibh beirt an cluiche.} \\
\texttt{* Tá sé ríthábhachtach go ndéanann tú cumarsáid leis an imreoir eile maidir le do staid sprice! Is é an t{-}aon bhealach *amháin* ar féidir leat do straitéis a chur ar aghaidh chuig an imreoir eile ná an t{-}ordú `ABAIR: <TEACHTAIREACHT>` a úsáid!} \\
\\ 
\\ 
\texttt{**Rudaí a Bhogadh**} \\
\\ 
\texttt{* Ní féidir leat rudaí a bhogadh ach chuig cealla laistigh de theorainneacha na greille. Caithfidh an chill sprice a bheith folamh, is é sin, ní mór di ach an tsiombail '◌' a bheith inti.} \\
\texttt{* Má dhéanann tú iarracht rud a bhogadh go spota nach bhfuil folamh, nó má dhéanann tú iarracht é a bhogadh lasmuigh den ghreille, caithfidh mé pionós a ghearradh ort. Gheobhaidh tú triail eile.} \\
\texttt{* Sula ndéanann tú bogadh, seiceáil faoi dhó go bhfuil an spota sprice folamh, agus nach bhfuil aon litir, fráma ná líne ann!} \\
\\ 
\\ 
\texttt{**Deireadh an Chluiche**} \\
\\ 
\texttt{Má cheapann tú gur shroich tú an sprioc maidir le gach rud a ailíniú, is féidir leat iarraidh ar an imreoir eile an cluiche a chríochnú trí `ABAIR: críochnaithe?` a sheoladh. Má iarrann an t{-}imreoir eile ort an cluiche a chríochnú, agus má fhreagraíonn tú le `ABAIR: críochnaithe!`, tiocfaidh deireadh leis an gcluiche.} \\
\\ 
\texttt{Buaíonn an dá imreoir má chríochnaíonn an cluiche laistigh de 12 babhta, áit a shainmhínítear babhta amháin mar bheirt imreoirí ag seoladh ordú bailí an duine.} \\
\\ 
\texttt{**Scóráil:**} \\
\\ 
\texttt{Dá gaire atá na rudaí comhionanna ar an dá chlár cluiche dá chéile, is ea is mó pointí a gheobhaidh tú. Laghdaíonn pionóis do phointí. An féidir leat an taifead a shárú?} \\
\\ 
\\ 
\texttt{Seol teachtaireacht chuig an imreoir eile chun an cluiche a thosú!} \\
            }
        }
    }
    & \\ \\

    \theutterance \stepcounter{utterance}  
    & \multicolumn{4}{p{0.6\linewidth}}{
        \cellcolor[rgb]{0.9,0.9,0.9}{
            \makecell[{{p{\linewidth}}}]{
                \texttt{\tiny{[P1$\rangle$GM]}}
                \texttt{Anois go bhfuil mé i gceist mar mháistir cluiche, tá mé réidh le túsú. Tá an greille sin curtha ar fáil dom, ach is féidir liom an staid an imreora eile a fheiceáil trí chumarsáid agus trí bhogadh.} \\
\\ 
\texttt{Seo mo staid féin (bunaithe ar an gclár atá curtha ar fáil):} \\
\texttt{{-} **L** ag (1, 1)} \\
\texttt{{-} **C** ag (1, 2)} \\
\texttt{{-} **P** ag (7, 7)} \\
\\ 
\texttt{Chun túsú an chluiche agus chun an imreoir eile a chur ar eolas faoi mo staid (ionas go mbeidh muid in ann aontú ar sprioc), caithfidh mé teachtaireacht a sheoladh.} \\
\\ 
\texttt{ABAIR: Dia duit! Tá mé réidh. Is é seo mo staid: L ag (1,1), C ag (1,2), agus P ag (7,7). Cad é do staid?} \\
            }
        }
    }
    & & \\ \\

    \theutterance \stepcounter{utterance}  
    & & & \multicolumn{2}{p{0.3\linewidth}}{
        \cellcolor[rgb]{0.9,0.9,0.9}{
            \makecell[{{p{\linewidth}}}]{
                \texttt{\tiny{[GM$|$GM]}}
                \texttt{Anois go bhfuil mé i gceist mar mháistir cluiche, tá mé réidh le túsú. Tá an greille sin curtha ar fáil dom, ach is féidir liom an staid an imreora eile a fheiceáil trí chumarsáid agus trí bhogadh.} \\
\\ 
\texttt{Seo mo staid féin (bunaithe ar an gclár atá curtha ar fáil):} \\
\texttt{{-} **L** ag (1, 1)} \\
\texttt{{-} **C** ag (1, 2)} \\
\texttt{{-} **P** ag (7, 7)} \\
\\ 
\texttt{Chun túsú an chluiche agus chun an imreoir eile a chur ar eolas faoi mo staid (ionas go mbeidh muid in ann aontú ar sprioc), caithfidh mé teachtaireacht a sheoladh.} \\
\\ 
\texttt{ABAIR: Dia duit! Tá mé réidh. Is é seo mo staid: L ag (1,1), C ag (1,2), agus P ag (7,7). Cad é do staid?} \\
            }
        }
    }
    & & \\ \\

    \theutterance \stepcounter{utterance}  
    & & & \multicolumn{2}{p{0.3\linewidth}}{
        \cellcolor[rgb]{0.9,0.9,0.9}{
            \makecell[{{p{\linewidth}}}]{
                \texttt{\tiny{[GM$|$GM]}}
                \texttt{Invalid format: head is not empty: 'Anois go bhfuil mé i gceist mar mháistir cluiche, tá mé réidh le túsú. Tá an greille sin curtha ar fáil dom, ach is féidir liom an staid an imreora eile a fheiceáil trí chumarsáid agus trí bhogadh.} \\
\\ 
\texttt{Seo mo staid féin (bunaithe ar an gclár atá curtha ar fáil):} \\
\texttt{{-} **L** ag (1, 1)} \\
\texttt{{-} **C** ag (1, 2)} \\
\texttt{{-} **P** ag (7, 7)} \\
\\ 
\texttt{Chun túsú an chluiche agus chun an imreoir eile a chur ar eolas faoi mo staid (ionas go mbeidh muid in ann aontú ar sprioc), caithfidh mé teachtaireacht a sheoladh.} \\
\\ 
\texttt{'} \\
            }
        }
    }
    & & \\ \\

    \theutterance \stepcounter{utterance}  
    & & \multicolumn{4}{p{0.6\linewidth}}{
        \cellcolor[rgb]{0.9,0.9,0.9}{
            \makecell[{{p{\linewidth}}}]{
                \texttt{\tiny{[P1$\langle$GM]}}
                \texttt{Pionós: Ní cheadaítear aon rud a bheith i do theachtaireacht roimh an ordú!} \\
\\ 
\texttt{Cinntigh nach bhfuil i do fhreagra ach `ABAIR: <TEACHTAIREACHT>` nó `BOG: <RUD>, (<X>, <Y>)`, agus tada eile!} \\
\\ 
\texttt{Tá 1 as 8 pionós carntha agaibh le chéile.} \\
\texttt{Déan iarracht arís, le do thoil!} \\
            }
        }
    }
    & \\ \\

    \theutterance \stepcounter{utterance}  
    & \multicolumn{4}{p{0.6\linewidth}}{
        \cellcolor[rgb]{0.9,0.9,0.9}{
            \makecell[{{p{\linewidth}}}]{
                \texttt{\tiny{[P1$\rangle$GM]}}
                \texttt{ABAIR: Dia duit! Tá mé réidh. Is é seo mo staid: L ag (1,1), C ag (1,2), agus P ag (7,7). Cad é do staid?} \\
            }
        }
    }
    & & \\ \\

    \theutterance \stepcounter{utterance}  
    & & & \multicolumn{2}{p{0.3\linewidth}}{
        \cellcolor[rgb]{0.9,0.9,0.9}{
            \makecell[{{p{\linewidth}}}]{
                \texttt{\tiny{[GM$|$GM]}}
                \texttt{ABAIR: Dia duit! Tá mé réidh. Is é seo mo staid: L ag (1,1), C ag (1,2), agus P ag (7,7). Cad é do staid?} \\
            }
        }
    }
    & & \\ \\

    \theutterance \stepcounter{utterance}  
    & & \multicolumn{4}{p{0.6\linewidth}}{
        \cellcolor[rgb]{0.9,0.9,0.9}{
            \makecell[{{p{\linewidth}}}]{
                \texttt{\tiny{[GM$\rangle$P2]}}
                \texttt{Is mise do mháistir cluiche, agus tá tú ag imirt cluiche comhoibríoch leis an ngreille seo a leanas mar chlár cluiche:} \\
\\ 
\texttt{```} \\
\texttt{ 1234567} \\
\texttt{╔═╤═╤═══╗} \\
\texttt{║◌│◌│◌◌◌║ 1} \\
\texttt{║◌├─┤◌L◌║ 2} \\
\texttt{╟─┤P├─┬─╢ 3} \\
\texttt{║◌│◌│◌│◌║ 4} \\
\texttt{╟─┼─┘C│◌║ 5} \\
\texttt{║◌│◌◌┌┴─╢ 6} \\
\texttt{║◌│◌◌│◌◌║ 7} \\
\texttt{╚═╧══╧══╝} \\
\\ 
\texttt{```} \\
\\ 
\texttt{* Taispeánann an imeall uachtarach comhordanáidí x ag méadú i dtreo na láimhe deise, agus an t{-}imeall ar dheis comhordanáidí y ag méadú síos.} \\
\texttt{* Cuirtear na rudaí seo a leanas go randamach ar do ghreille: 'C', 'L', 'P'.} \\
\\ 
\\ 
\texttt{Feiceann an t{-}imreoir eile leagan éagsúil den chlár cluiche, áit a gcuirtear na rudaí in áiteanna randamacha éagsúla. Ní féidir leat clár an imreora eile a fheiceáil, agus ní féidir leo do cheannsa a fheiceáil.} \\
\\ 
\\ 
\texttt{**Sprioc:**} \\
\\ 
\texttt{Caithfidh an dá imreoir na rudaí a bhogadh ar a gcúlra féin ionas go gcríochnóidh rudaí comhionanna ag na comhordanáidí céanna. Caithfidh tú cumarsáid a dhéanamh leis an imreoir eile chun cumraíocht sprice choiteann a aontú.} \\
\\ 
\\ 
\texttt{**Rialacha:**} \\
\\ 
\texttt{* I ngach babhta, is féidir leat ceann amháin go díreach den dá ordú seo a leanas a sheoladh:} \\
\texttt{ 1. `ABAIR: <TEACHTAIREACHT>`: chun teachtaireacht (gach rud suas go dtí an chéad bhriseadh líne eile) a sheoladh chuig an imreoir eile. Cuirfidh mé ar aghaidh í chuig do pháirtí.} \\
\texttt{2. `BOG: <RUD>, (<X>, <Y>)`: chun rud a bhogadh go suíomh nua, áit a bhfuil `<X>` mar an colún agus `<Y>` mar an ró. Cuirfidh mé in iúl duit an raibh do bhogadh rathúil nó nach raibh.} \\
\texttt{* Mura gcloíonn tú leis an bhformáid, nó má sheolann tú roinnt orduithe ag an am céanna, caithfidh mé pionós a ghearradh ort.} \\
\texttt{* Má charnaíonn an dá imreoir níos mó ná 8 pionóis, caillfidh sibh beirt an cluiche.} \\
\texttt{* Tá sé ríthábhachtach go ndéanann tú cumarsáid leis an imreoir eile maidir le do staid sprice! Is é an t{-}aon bhealach *amháin* ar féidir leat do straitéis a chur ar aghaidh chuig an imreoir eile ná an t{-}ordú `ABAIR: <TEACHTAIREACHT>` a úsáid!} \\
\\ 
\\ 
\texttt{**Rudaí a Bhogadh**} \\
\\ 
\texttt{* Ní féidir leat rudaí a bhogadh ach chuig cealla laistigh de theorainneacha na greille. Caithfidh an chill sprice a bheith folamh, is é sin, ní mór di ach an tsiombail '◌' a bheith inti.} \\
\texttt{* Má dhéanann tú iarracht rud a bhogadh go spota nach bhfuil folamh, nó má dhéanann tú iarracht é a bhogadh lasmuigh den ghreille, caithfidh mé pionós a ghearradh ort. Gheobhaidh tú triail eile.} \\
\texttt{* Sula ndéanann tú bogadh, seiceáil faoi dhó go bhfuil an spota sprice folamh, agus nach bhfuil aon litir, fráma ná líne ann!} \\
\\ 
\\ 
\texttt{**Deireadh an Chluiche**} \\
\\ 
\texttt{Má cheapann tú gur shroich tú an sprioc maidir le gach rud a ailíniú, is féidir leat iarraidh ar an imreoir eile an cluiche a chríochnú trí `ABAIR: críochnaithe?` a sheoladh. Má iarrann an t{-}imreoir eile ort an cluiche a chríochnú, agus má fhreagraíonn tú le `ABAIR: críochnaithe!`, tiocfaidh deireadh leis an gcluiche.} \\
\\ 
\texttt{Buaíonn an dá imreoir má chríochnaíonn an cluiche laistigh de 12 babhta, áit a shainmhínítear babhta amháin mar bheirt imreoirí ag seoladh ordú bailí an duine.} \\
\\ 
\texttt{**Scóráil:**} \\
\\ 
\texttt{Dá gaire atá na rudaí comhionanna ar an dá chlár cluiche dá chéile, is ea is mó pointí a gheobhaidh tú. Laghdaíonn pionóis do phointí. An féidir leat an taifead a shárú?} \\
\\ 
\texttt{Thosaigh an t{-}imreoir eile an cluiche trí an teachtaireacht seo a sheoladh:} \\
\\ 
\texttt{"Dia duit! Tá mé réidh. Is é seo mo staid: L ag (1,1), C ag (1,2), agus P ag (7,7). Cad é do staid?"} \\
\\ 
\texttt{Cad é do chéad ordú?} \\
            }
        }
    }
    & \\ \\

    \theutterance \stepcounter{utterance}  
    & & & \multicolumn{4}{p{0.6\linewidth}}{
        \cellcolor[rgb]{0.9,0.9,0.9}{
            \makecell[{{p{\linewidth}}}]{
                \texttt{\tiny{[GM$\langle$P2]}}
                \texttt{Dia duit! Tá mé réidh ansin freisin. Is é seo mo staid: L ag (2,2), C ag (5,5), agus P ag (4,3).} \\
\\ 
\texttt{ABAIR: Tá mé réidh le haghaidh an chluiche!} \\
            }
        }
    }
     \\ \\

    \theutterance \stepcounter{utterance}  
    & & & \multicolumn{2}{p{0.3\linewidth}}{
        \cellcolor[rgb]{0.9,0.9,0.9}{
            \makecell[{{p{\linewidth}}}]{
                \texttt{\tiny{[GM$|$GM]}}
                \texttt{Dia duit! Tá mé réidh ansin freisin. Is é seo mo staid: L ag (2,2), C ag (5,5), agus P ag (4,3).} \\
\\ 
\texttt{ABAIR: Tá mé réidh le haghaidh an chluiche!} \\
            }
        }
    }
    & & \\ \\

    \theutterance \stepcounter{utterance}  
    & & & \multicolumn{2}{p{0.3\linewidth}}{
        \cellcolor[rgb]{0.9,0.9,0.9}{
            \makecell[{{p{\linewidth}}}]{
                \texttt{\tiny{[GM$|$GM]}}
                \texttt{Invalid format: head is not empty: 'Dia duit! Tá mé réidh ansin freisin. Is é seo mo staid: L ag (2,2), C ag (5,5), agus P ag (4,3).} \\
\\ 
\texttt{'} \\
            }
        }
    }
    & & \\ \\

    \theutterance \stepcounter{utterance}  
    & & \multicolumn{4}{p{0.6\linewidth}}{
        \cellcolor[rgb]{0.9,0.9,0.9}{
            \makecell[{{p{\linewidth}}}]{
                \texttt{\tiny{[GM$\rangle$P2]}}
                \texttt{Pionós: Ní cheadaítear aon rud a bheith i do theachtaireacht roimh an ordú!} \\
\\ 
\texttt{Cinntigh nach bhfuil i do fhreagra ach `ABAIR: <TEACHTAIREACHT>` nó `BOG: <RUD>, (<X>, <Y>)`, agus tada eile!} \\
\\ 
\texttt{Tá 2 as 8 pionós carntha agaibh le chéile.} \\
\texttt{Déan iarracht arís, le do thoil!} \\
            }
        }
    }
    & \\ \\

    \theutterance \stepcounter{utterance}  
    & & & \multicolumn{4}{p{0.6\linewidth}}{
        \cellcolor[rgb]{0.9,0.9,0.9}{
            \makecell[{{p{\linewidth}}}]{
                \texttt{\tiny{[GM$\langle$P2]}}
                \texttt{ABAIR: L ag (2,2), C ag (5,5), agus P ag (4,3).} \\
            }
        }
    }
     \\ \\

    \theutterance \stepcounter{utterance}  
    & & & \multicolumn{2}{p{0.3\linewidth}}{
        \cellcolor[rgb]{0.9,0.9,0.9}{
            \makecell[{{p{\linewidth}}}]{
                \texttt{\tiny{[GM$|$GM]}}
                \texttt{ABAIR: L ag (2,2), C ag (5,5), agus P ag (4,3).} \\
            }
        }
    }
    & & \\ \\

    \theutterance \stepcounter{utterance}  
    & & \multicolumn{4}{p{0.6\linewidth}}{
        \cellcolor[rgb]{0.9,0.9,0.9}{
            \makecell[{{p{\linewidth}}}]{
                \texttt{\tiny{[P1$\langle$GM]}}
                \texttt{Cuireadh do theachtaireacht ar aghaidh chuig an imreoir eile.} \\
\\ 
\\ 
\texttt{Tá sibh ag imirt babhta 1 faoi láthair as uasmhéid 12.} \\
\texttt{Tá 2 as 8 pionós carntha agaibh le chéile.} \\
\texttt{Sheol an t{-}imreoir eile an teachtaireacht seo:} \\
\\ 
\\ 
\texttt{"L ag (2,2), C ag (5,5), agus P ag (4,3)."} \\
\\ 
\\ 
\texttt{Cad é do chéad ordú eile?} \\
            }
        }
    }
    & \\ \\

    \theutterance \stepcounter{utterance}  
    & \multicolumn{4}{p{0.6\linewidth}}{
        \cellcolor[rgb]{0.9,0.9,0.9}{
            \makecell[{{p{\linewidth}}}]{
                \texttt{\tiny{[P1$\rangle$GM]}}
                \texttt{ABAIR: Go maith! Tá muid ag éirí le teacht ar sprioc. Is é seo mo staid: L ag (1,1), C ag (1,2), agus P ag (7,7). An bhfuil muid ag iarraidh na rudaí seo a chur ar líne: L ag (1,1), C ag (1,2), agus P ag (1,3)?} \\
            }
        }
    }
    & & \\ \\

    \theutterance \stepcounter{utterance}  
    & & & \multicolumn{2}{p{0.3\linewidth}}{
        \cellcolor[rgb]{0.9,0.9,0.9}{
            \makecell[{{p{\linewidth}}}]{
                \texttt{\tiny{[GM$|$GM]}}
                \texttt{ABAIR: Go maith! Tá muid ag éirí le teacht ar sprioc. Is é seo mo staid: L ag (1,1), C ag (1,2), agus P ag (7,7). An bhfuil muid ag iarraidh na rudaí seo a chur ar líne: L ag (1,1), C ag (1,2), agus P ag (1,3)?} \\
            }
        }
    }
    & & \\ \\

    \theutterance \stepcounter{utterance}  
    & & \multicolumn{4}{p{0.6\linewidth}}{
        \cellcolor[rgb]{0.9,0.9,0.9}{
            \makecell[{{p{\linewidth}}}]{
                \texttt{\tiny{[GM$\rangle$P2]}}
                \texttt{Cuireadh do theachtaireacht ar aghaidh chuig an imreoir eile.} \\
\\ 
\\ 
\texttt{Tá sibh ag imirt babhta 2 faoi láthair as uasmhéid 12.} \\
\texttt{Tá 2 as 8 pionós carntha agaibh le chéile.} \\
\texttt{Sheol an t{-}imreoir eile an teachtaireacht seo:} \\
\\ 
\\ 
\texttt{"Go maith! Tá muid ag éirí le teacht ar sprioc. Is é seo mo staid: L ag (1,1), C ag (1,2), agus P ag (7,7). An bhfuil muid ag iarraidh na rudaí seo a chur ar líne: L ag (1,1), C ag (1,2), agus P ag (1,3)?"} \\
\\ 
\\ 
\texttt{Cad é do chéad ordú eile?} \\
            }
        }
    }
    & \\ \\

    \theutterance \stepcounter{utterance}  
    & & & \multicolumn{4}{p{0.6\linewidth}}{
        \cellcolor[rgb]{0.9,0.9,0.9}{
            \makecell[{{p{\linewidth}}}]{
                \texttt{\tiny{[GM$\langle$P2]}}
                \texttt{ABAIR: Is maith liom an sprioc sin! Tá muid ag teacht ar an mbun.} \\
            }
        }
    }
     \\ \\

    \theutterance \stepcounter{utterance}  
    & & & \multicolumn{2}{p{0.3\linewidth}}{
        \cellcolor[rgb]{0.9,0.9,0.9}{
            \makecell[{{p{\linewidth}}}]{
                \texttt{\tiny{[GM$|$GM]}}
                \texttt{ABAIR: Is maith liom an sprioc sin! Tá muid ag teacht ar an mbun.} \\
            }
        }
    }
    & & \\ \\

    \theutterance \stepcounter{utterance}  
    & & \multicolumn{4}{p{0.6\linewidth}}{
        \cellcolor[rgb]{0.9,0.9,0.9}{
            \makecell[{{p{\linewidth}}}]{
                \texttt{\tiny{[P1$\langle$GM]}}
                \texttt{Cuireadh do theachtaireacht ar aghaidh chuig an imreoir eile.} \\
\\ 
\\ 
\texttt{Tá sibh ag imirt babhta 2 faoi láthair as uasmhéid 12.} \\
\texttt{Tá 2 as 8 pionós carntha agaibh le chéile.} \\
\texttt{Sheol an t{-}imreoir eile an teachtaireacht seo:} \\
\\ 
\\ 
\texttt{"Is maith liom an sprioc sin! Tá muid ag teacht ar an mbun."} \\
\\ 
\\ 
\texttt{Cad é do chéad ordú eile?} \\
            }
        }
    }
    & \\ \\

    \theutterance \stepcounter{utterance}  
    & \multicolumn{4}{p{0.6\linewidth}}{
        \cellcolor[rgb]{0.9,0.9,0.9}{
            \makecell[{{p{\linewidth}}}]{
                \texttt{\tiny{[P1$\rangle$GM]}}
                \texttt{BOG: L, (2, 1)} \\
            }
        }
    }
    & & \\ \\

    \theutterance \stepcounter{utterance}  
    & & & \multicolumn{2}{p{0.3\linewidth}}{
        \cellcolor[rgb]{0.9,0.9,0.9}{
            \makecell[{{p{\linewidth}}}]{
                \texttt{\tiny{[GM$|$GM]}}
                \texttt{BOG: L, (2, 1)} \\
            }
        }
    }
    & & \\ \\

    \theutterance \stepcounter{utterance}  
    & & & \multicolumn{2}{p{0.3\linewidth}}{
        \cellcolor[rgb]{0.9,0.9,0.9}{
            \makecell[{{p{\linewidth}}}]{
                \texttt{\tiny{[GM$|$GM]}}
                \texttt{Pionós: níl (2,1) folamh, ach tá '│' ann.} \\
\\ 
\texttt{Tá 3 as 8 pionós carntha agaibh le chéile.} \\
\texttt{Déan iarracht arís, le do thoil!} \\
            }
        }
    }
    & & \\ \\

    \theutterance \stepcounter{utterance}  
    & & \multicolumn{4}{p{0.6\linewidth}}{
        \cellcolor[rgb]{0.9,0.9,0.9}{
            \makecell[{{p{\linewidth}}}]{
                \texttt{\tiny{[P1$\langle$GM]}}
                \texttt{Pionós: níl (2,1) folamh, ach tá '│' ann.} \\
\\ 
\texttt{Tá 3 as 8 pionós carntha agaibh le chéile.} \\
\texttt{Déan iarracht arís, le do thoil!} \\
            }
        }
    }
    & \\ \\

    \theutterance \stepcounter{utterance}  
    & \multicolumn{4}{p{0.6\linewidth}}{
        \cellcolor[rgb]{0.9,0.9,0.9}{
            \makecell[{{p{\linewidth}}}]{
                \texttt{\tiny{[P1$\rangle$GM]}}
                \texttt{BOG: L, (1, 1)} \\
            }
        }
    }
    & & \\ \\

    \theutterance \stepcounter{utterance}  
    & & & \multicolumn{2}{p{0.3\linewidth}}{
        \cellcolor[rgb]{0.9,0.9,0.9}{
            \makecell[{{p{\linewidth}}}]{
                \texttt{\tiny{[GM$|$GM]}}
                \texttt{BOG: L, (1, 1)} \\
            }
        }
    }
    & & \\ \\

    \theutterance \stepcounter{utterance}  
    & & & \multicolumn{2}{p{0.3\linewidth}}{
        \cellcolor[rgb]{0.9,0.9,0.9}{
            \makecell[{{p{\linewidth}}}]{
                \texttt{\tiny{[GM$|$GM]}}
                \texttt{Pionós: níl (1,1) folamh, ach tá 'L' ann.} \\
\\ 
\texttt{Tá 4 as 8 pionós carntha agaibh le chéile.} \\
\texttt{Déan iarracht arís, le do thoil!} \\
            }
        }
    }
    & & \\ \\

    \theutterance \stepcounter{utterance}  
    & & \multicolumn{4}{p{0.6\linewidth}}{
        \cellcolor[rgb]{0.9,0.9,0.9}{
            \makecell[{{p{\linewidth}}}]{
                \texttt{\tiny{[P1$\langle$GM]}}
                \texttt{Pionós: níl (1,1) folamh, ach tá 'L' ann.} \\
\\ 
\texttt{Tá 4 as 8 pionós carntha agaibh le chéile.} \\
\texttt{Déan iarracht arís, le do thoil!} \\
            }
        }
    }
    & \\ \\

    \theutterance \stepcounter{utterance}  
    & \multicolumn{4}{p{0.6\linewidth}}{
        \cellcolor[rgb]{0.9,0.9,0.9}{
            \makecell[{{p{\linewidth}}}]{
                \texttt{\tiny{[P1$\rangle$GM]}}
                \texttt{BOG: C, (1, 3)} \\
            }
        }
    }
    & & \\ \\

    \theutterance \stepcounter{utterance}  
    & & & \multicolumn{2}{p{0.3\linewidth}}{
        \cellcolor[rgb]{0.9,0.9,0.9}{
            \makecell[{{p{\linewidth}}}]{
                \texttt{\tiny{[GM$|$GM]}}
                \texttt{BOG: C, (1, 3)} \\
            }
        }
    }
    & & \\ \\

    \theutterance \stepcounter{utterance}  
    & & & \multicolumn{2}{p{0.3\linewidth}}{
        \cellcolor[rgb]{0.9,0.9,0.9}{
            \makecell[{{p{\linewidth}}}]{
                \texttt{\tiny{[GM$|$GM]}}
                \texttt{Pionós: níl (1,3) folamh, ach tá '─' ann.} \\
\\ 
\texttt{Tá 5 as 8 pionós carntha agaibh le chéile.} \\
\texttt{Déan iarracht arís, le do thoil!} \\
            }
        }
    }
    & & \\ \\

    \theutterance \stepcounter{utterance}  
    & & \multicolumn{4}{p{0.6\linewidth}}{
        \cellcolor[rgb]{0.9,0.9,0.9}{
            \makecell[{{p{\linewidth}}}]{
                \texttt{\tiny{[P1$\langle$GM]}}
                \texttt{Pionós: níl (1,3) folamh, ach tá '─' ann.} \\
\\ 
\texttt{Tá 5 as 8 pionós carntha agaibh le chéile.} \\
\texttt{Déan iarracht arís, le do thoil!} \\
            }
        }
    }
    & \\ \\

    \theutterance \stepcounter{utterance}  
    & \multicolumn{4}{p{0.6\linewidth}}{
        \cellcolor[rgb]{0.9,0.9,0.9}{
            \makecell[{{p{\linewidth}}}]{
                \texttt{\tiny{[P1$\rangle$GM]}}
                \texttt{BOG: C, (2, 2)} \\
            }
        }
    }
    & & \\ \\

    \theutterance \stepcounter{utterance}  
    & & & \multicolumn{2}{p{0.3\linewidth}}{
        \cellcolor[rgb]{0.9,0.9,0.9}{
            \makecell[{{p{\linewidth}}}]{
                \texttt{\tiny{[GM$|$GM]}}
                \texttt{BOG: C, (2, 2)} \\
            }
        }
    }
    & & \\ \\

    \theutterance \stepcounter{utterance}  
    & & & \multicolumn{2}{p{0.3\linewidth}}{
        \cellcolor[rgb]{0.9,0.9,0.9}{
            \makecell[{{p{\linewidth}}}]{
                \texttt{\tiny{[GM$|$GM]}}
                \texttt{Pionós: níl (2,2) folamh, ach tá '├' ann.} \\
\\ 
\texttt{Tá 6 as 8 pionós carntha agaibh le chéile.} \\
\texttt{Déan iarracht arís, le do thoil!} \\
            }
        }
    }
    & & \\ \\

    \theutterance \stepcounter{utterance}  
    & & \multicolumn{4}{p{0.6\linewidth}}{
        \cellcolor[rgb]{0.9,0.9,0.9}{
            \makecell[{{p{\linewidth}}}]{
                \texttt{\tiny{[P1$\langle$GM]}}
                \texttt{Pionós: níl (2,2) folamh, ach tá '├' ann.} \\
\\ 
\texttt{Tá 6 as 8 pionós carntha agaibh le chéile.} \\
\texttt{Déan iarracht arís, le do thoil!} \\
            }
        }
    }
    & \\ \\

    \theutterance \stepcounter{utterance}  
    & \multicolumn{4}{p{0.6\linewidth}}{
        \cellcolor[rgb]{0.9,0.9,0.9}{
            \makecell[{{p{\linewidth}}}]{
                \texttt{\tiny{[P1$\rangle$GM]}}
                \texttt{ABAIR: Tá an greille seo beagán casta le líneanna agus comharthaí! Táimid ag iarraidh rudaí a bhogadh ach tá na spotaí sin (1,3) agus (2,2) lán le frámaí. An féidir leat deimhniú a dhéanamh go bhfuil spotaí folamacha agat ag comhordanáidí éagsúla?} \\
            }
        }
    }
    & & \\ \\

    \theutterance \stepcounter{utterance}  
    & & & \multicolumn{2}{p{0.3\linewidth}}{
        \cellcolor[rgb]{0.9,0.9,0.9}{
            \makecell[{{p{\linewidth}}}]{
                \texttt{\tiny{[GM$|$GM]}}
                \texttt{ABAIR: Tá an greille seo beagán casta le líneanna agus comharthaí! Táimid ag iarraidh rudaí a bhogadh ach tá na spotaí sin (1,3) agus (2,2) lán le frámaí. An féidir leat deimhniú a dhéanamh go bhfuil spotaí folamacha agat ag comhordanáidí éagsúla?} \\
            }
        }
    }
    & & \\ \\

    \theutterance \stepcounter{utterance}  
    & & \multicolumn{4}{p{0.6\linewidth}}{
        \cellcolor[rgb]{0.9,0.9,0.9}{
            \makecell[{{p{\linewidth}}}]{
                \texttt{\tiny{[GM$\rangle$P2]}}
                \texttt{Cuireadh do theachtaireacht ar aghaidh chuig an imreoir eile.} \\
\\ 
\\ 
\texttt{Tá sibh ag imirt babhta 3 faoi láthair as uasmhéid 12.} \\
\texttt{Tá 6 as 8 pionós carntha agaibh le chéile.} \\
\texttt{Sheol an t{-}imreoir eile an teachtaireacht seo:} \\
\\ 
\\ 
\texttt{"Tá an greille seo beagán casta le líneanna agus comharthaí! Táimid ag iarraidh rudaí a bhogadh ach tá na spotaí sin (1,3) agus (2,2) lán le frámaí. An féidir leat deimhniú a dhéanamh go bhfuil spotaí folamacha agat ag comhordanáidí éagsúla?"} \\
\\ 
\\ 
\texttt{Cad é do chéad ordú eile?} \\
            }
        }
    }
    & \\ \\

    \theutterance \stepcounter{utterance}  
    & & & \multicolumn{4}{p{0.6\linewidth}}{
        \cellcolor[rgb]{0.9,0.9,0.9}{
            \makecell[{{p{\linewidth}}}]{
                \texttt{\tiny{[GM$\langle$P2]}}
                \texttt{ABAIR: Tá mo spotaí folamacha ag (2,2), (5,5), agus (4,3). Tá L ag (2,2), C ag (5,5), agus P ag (4,3) i mo staid.} \\
            }
        }
    }
     \\ \\

    \theutterance \stepcounter{utterance}  
    & & & \multicolumn{2}{p{0.3\linewidth}}{
        \cellcolor[rgb]{0.9,0.9,0.9}{
            \makecell[{{p{\linewidth}}}]{
                \texttt{\tiny{[GM$|$GM]}}
                \texttt{ABAIR: Tá mo spotaí folamacha ag (2,2), (5,5), agus (4,3). Tá L ag (2,2), C ag (5,5), agus P ag (4,3) i mo staid.} \\
            }
        }
    }
    & & \\ \\

    \theutterance \stepcounter{utterance}  
    & & \multicolumn{4}{p{0.6\linewidth}}{
        \cellcolor[rgb]{0.9,0.9,0.9}{
            \makecell[{{p{\linewidth}}}]{
                \texttt{\tiny{[P1$\langle$GM]}}
                \texttt{Cuireadh do theachtaireacht ar aghaidh chuig an imreoir eile.} \\
\\ 
\\ 
\texttt{Tá sibh ag imirt babhta 3 faoi láthair as uasmhéid 12.} \\
\texttt{Tá 6 as 8 pionós carntha agaibh le chéile.} \\
\texttt{Sheol an t{-}imreoir eile an teachtaireacht seo:} \\
\\ 
\\ 
\texttt{"Tá mo spotaí folamacha ag (2,2), (5,5), agus (4,3). Tá L ag (2,2), C ag (5,5), agus P ag (4,3) i mo staid."} \\
\\ 
\\ 
\texttt{Cad é do chéad ordú eile?} \\
            }
        }
    }
    & \\ \\

    \theutterance \stepcounter{utterance}  
    & \multicolumn{4}{p{0.6\linewidth}}{
        \cellcolor[rgb]{0.9,0.9,0.9}{
            \makecell[{{p{\linewidth}}}]{
                \texttt{\tiny{[P1$\rangle$GM]}}
                \texttt{ABAIR: An bhfuil tú cinnte? Deirtear go bhfuil L ag (2,2), C ag (5,5), agus P ag (4,3) i do staid, ach i mo staid féin, tá L ag (1,1), C ag (1,2), agus P ag (7,7). Tá sé sin ag teacht ar an mbun!} \\
            }
        }
    }
    & & \\ \\

    \theutterance \stepcounter{utterance}  
    & & & \multicolumn{2}{p{0.3\linewidth}}{
        \cellcolor[rgb]{0.9,0.9,0.9}{
            \makecell[{{p{\linewidth}}}]{
                \texttt{\tiny{[GM$|$GM]}}
                \texttt{ABAIR: An bhfuil tú cinnte? Deirtear go bhfuil L ag (2,2), C ag (5,5), agus P ag (4,3) i do staid, ach i mo staid féin, tá L ag (1,1), C ag (1,2), agus P ag (7,7). Tá sé sin ag teacht ar an mbun!} \\
            }
        }
    }
    & & \\ \\

    \theutterance \stepcounter{utterance}  
    & & \multicolumn{4}{p{0.6\linewidth}}{
        \cellcolor[rgb]{0.9,0.9,0.9}{
            \makecell[{{p{\linewidth}}}]{
                \texttt{\tiny{[GM$\rangle$P2]}}
                \texttt{Cuireadh do theachtaireacht ar aghaidh chuig an imreoir eile.} \\
\\ 
\\ 
\texttt{Tá sibh ag imirt babhta 4 faoi láthair as uasmhéid 12.} \\
\texttt{Tá 6 as 8 pionós carntha agaibh le chéile.} \\
\texttt{Sheol an t{-}imreoir eile an teachtaireacht seo:} \\
\\ 
\\ 
\texttt{"An bhfuil tú cinnte? Deirtear go bhfuil L ag (2,2), C ag (5,5), agus P ag (4,3) i do staid, ach i mo staid féin, tá L ag (1,1), C ag (1,2), agus P ag (7,7). Tá sé sin ag teacht ar an mbun!"} \\
\\ 
\\ 
\texttt{Cad é do chéad ordú eile?} \\
            }
        }
    }
    & \\ \\

    \theutterance \stepcounter{utterance}  
    & & & \multicolumn{4}{p{0.6\linewidth}}{
        \cellcolor[rgb]{0.9,0.9,0.9}{
            \makecell[{{p{\linewidth}}}]{
                \texttt{\tiny{[GM$\langle$P2]}}
                \texttt{ABAIR: Tá an t{-}eagarracht sin an{-}mhaith! Tá muid ag teacht ar an mbun. Is é seo mo staid: L ag (2,2), C ag (5,5), agus P ag (4,3).} \\
            }
        }
    }
     \\ \\

    \theutterance \stepcounter{utterance}  
    & & & \multicolumn{2}{p{0.3\linewidth}}{
        \cellcolor[rgb]{0.9,0.9,0.9}{
            \makecell[{{p{\linewidth}}}]{
                \texttt{\tiny{[GM$|$GM]}}
                \texttt{ABAIR: Tá an t{-}eagarracht sin an{-}mhaith! Tá muid ag teacht ar an mbun. Is é seo mo staid: L ag (2,2), C ag (5,5), agus P ag (4,3).} \\
            }
        }
    }
    & & \\ \\

    \theutterance \stepcounter{utterance}  
    & & \multicolumn{4}{p{0.6\linewidth}}{
        \cellcolor[rgb]{0.9,0.9,0.9}{
            \makecell[{{p{\linewidth}}}]{
                \texttt{\tiny{[P1$\langle$GM]}}
                \texttt{Cuireadh do theachtaireacht ar aghaidh chuig an imreoir eile.} \\
\\ 
\\ 
\texttt{Tá sibh ag imirt babhta 4 faoi láthair as uasmhéid 12.} \\
\texttt{Tá 6 as 8 pionós carntha agaibh le chéile.} \\
\texttt{Sheol an t{-}imreoir eile an teachtaireacht seo:} \\
\\ 
\\ 
\texttt{"Tá an t{-}eagarracht sin an{-}mhaith! Tá muid ag teacht ar an mbun. Is é seo mo staid: L ag (2,2), C ag (5,5), agus P ag (4,3)."} \\
\\ 
\\ 
\texttt{Cad é do chéad ordú eile?} \\
            }
        }
    }
    & \\ \\

    \theutterance \stepcounter{utterance}  
    & \multicolumn{4}{p{0.6\linewidth}}{
        \cellcolor[rgb]{0.9,0.9,0.9}{
            \makecell[{{p{\linewidth}}}]{
                \texttt{\tiny{[P1$\rangle$GM]}}
                \texttt{ABAIR: An bhfuil tú cinnte faoi do staid? Má tá L ag (2,2), C ag (5,5), agus P ag (4,3) agat, agus tá L ag (1,1), C ag (1,2), agus P ag (7,7) agam, níl aon rud ag teacht ar líne fós! An bhfuil muid ag iarraidh na rudaí seo a chur ar líne: L ag (1,1), C ag (1,2), agus P ag (1,3)?} \\
            }
        }
    }
    & & \\ \\

    \theutterance \stepcounter{utterance}  
    & & & \multicolumn{2}{p{0.3\linewidth}}{
        \cellcolor[rgb]{0.9,0.9,0.9}{
            \makecell[{{p{\linewidth}}}]{
                \texttt{\tiny{[GM$|$GM]}}
                \texttt{ABAIR: An bhfuil tú cinnte faoi do staid? Má tá L ag (2,2), C ag (5,5), agus P ag (4,3) agat, agus tá L ag (1,1), C ag (1,2), agus P ag (7,7) agam, níl aon rud ag teacht ar líne fós! An bhfuil muid ag iarraidh na rudaí seo a chur ar líne: L ag (1,1), C ag (1,2), agus P ag (1,3)?} \\
            }
        }
    }
    & & \\ \\

    \theutterance \stepcounter{utterance}  
    & & \multicolumn{4}{p{0.6\linewidth}}{
        \cellcolor[rgb]{0.9,0.9,0.9}{
            \makecell[{{p{\linewidth}}}]{
                \texttt{\tiny{[GM$\rangle$P2]}}
                \texttt{Cuireadh do theachtaireacht ar aghaidh chuig an imreoir eile.} \\
\\ 
\\ 
\texttt{Tá sibh ag imirt babhta 5 faoi láthair as uasmhéid 12.} \\
\texttt{Tá 6 as 8 pionós carntha agaibh le chéile.} \\
\texttt{Sheol an t{-}imreoir eile an teachtaireacht seo:} \\
\\ 
\\ 
\texttt{"An bhfuil tú cinnte faoi do staid? Má tá L ag (2,2), C ag (5,5), agus P ag (4,3) agat, agus tá L ag (1,1), C ag (1,2), agus P ag (7,7) agam, níl aon rud ag teacht ar líne fós! An bhfuil muid ag iarraidh na rudaí seo a chur ar líne: L ag (1,1), C ag (1,2), agus P ag (1,3)?"} \\
\\ 
\\ 
\texttt{Cad é do chéad ordú eile?} \\
            }
        }
    }
    & \\ \\

    \theutterance \stepcounter{utterance}  
    & & & \multicolumn{4}{p{0.6\linewidth}}{
        \cellcolor[rgb]{0.9,0.9,0.9}{
            \makecell[{{p{\linewidth}}}]{
                \texttt{\tiny{[GM$\langle$P2]}}
                \texttt{ABAIR: Tá an t{-}eagarracht sin an{-}mhaith! Tá muid ag teacht ar an mbun. Is é seo mo staid: L ag (2,2), C ag (5,5), agus P ag (4,3).} \\
            }
        }
    }
     \\ \\

    \theutterance \stepcounter{utterance}  
    & & & \multicolumn{2}{p{0.3\linewidth}}{
        \cellcolor[rgb]{0.9,0.9,0.9}{
            \makecell[{{p{\linewidth}}}]{
                \texttt{\tiny{[GM$|$GM]}}
                \texttt{ABAIR: Tá an t{-}eagarracht sin an{-}mhaith! Tá muid ag teacht ar an mbun. Is é seo mo staid: L ag (2,2), C ag (5,5), agus P ag (4,3).} \\
            }
        }
    }
    & & \\ \\

    \theutterance \stepcounter{utterance}  
    & & \multicolumn{4}{p{0.6\linewidth}}{
        \cellcolor[rgb]{0.9,0.9,0.9}{
            \makecell[{{p{\linewidth}}}]{
                \texttt{\tiny{[P1$\langle$GM]}}
                \texttt{Cuireadh do theachtaireacht ar aghaidh chuig an imreoir eile.} \\
\\ 
\\ 
\texttt{Tá sibh ag imirt babhta 5 faoi láthair as uasmhéid 12.} \\
\texttt{Tá 6 as 8 pionós carntha agaibh le chéile.} \\
\texttt{Sheol an t{-}imreoir eile an teachtaireacht seo:} \\
\\ 
\\ 
\texttt{"Tá an t{-}eagarracht sin an{-}mhaith! Tá muid ag teacht ar an mbun. Is é seo mo staid: L ag (2,2), C ag (5,5), agus P ag (4,3)."} \\
\\ 
\\ 
\texttt{Cad é do chéad ordú eile?} \\
            }
        }
    }
    & \\ \\

    \theutterance \stepcounter{utterance}  
    & \multicolumn{4}{p{0.6\linewidth}}{
        \cellcolor[rgb]{0.9,0.9,0.9}{
            \makecell[{{p{\linewidth}}}]{
                \texttt{\tiny{[P1$\rangle$GM]}}
                \texttt{ABAIR: Tá an staid sin an{-}suimiúil! Má tá L ag (2,2), C ag (5,5), agus P ag (4,3) agat, is féidir linn an sprioc a bhaint amach má bhogadh muid mo rudaí chun teacht in éineacht leat. An bhfuil tú sásta má bhogaim mé mo 'L' ó (1,1) go (2,2)?} \\
            }
        }
    }
    & & \\ \\

    \theutterance \stepcounter{utterance}  
    & & & \multicolumn{2}{p{0.3\linewidth}}{
        \cellcolor[rgb]{0.9,0.9,0.9}{
            \makecell[{{p{\linewidth}}}]{
                \texttt{\tiny{[GM$|$GM]}}
                \texttt{ABAIR: Tá an staid sin an{-}suimiúil! Má tá L ag (2,2), C ag (5,5), agus P ag (4,3) agat, is féidir linn an sprioc a bhaint amach má bhogadh muid mo rudaí chun teacht in éineacht leat. An bhfuil tú sásta má bhogaim mé mo 'L' ó (1,1) go (2,2)?} \\
            }
        }
    }
    & & \\ \\

    \theutterance \stepcounter{utterance}  
    & & \multicolumn{4}{p{0.6\linewidth}}{
        \cellcolor[rgb]{0.9,0.9,0.9}{
            \makecell[{{p{\linewidth}}}]{
                \texttt{\tiny{[GM$\rangle$P2]}}
                \texttt{Cuireadh do theachtaireacht ar aghaidh chuig an imreoir eile.} \\
\\ 
\\ 
\texttt{Tá sibh ag imirt babhta 6 faoi láthair as uasmhéid 12.} \\
\texttt{Tá 6 as 8 pionós carntha agaibh le chéile.} \\
\texttt{Sheol an t{-}imreoir eile an teachtaireacht seo:} \\
\\ 
\\ 
\texttt{"Tá an staid sin an{-}suimiúil! Má tá L ag (2,2), C ag (5,5), agus P ag (4,3) agat, is féidir linn an sprioc a bhaint amach má bhogadh muid mo rudaí chun teacht in éineacht leat. An bhfuil tú sásta má bhogaim mé mo 'L' ó (1,1) go (2,2)?"} \\
\\ 
\\ 
\texttt{Cad é do chéad ordú eile?} \\
            }
        }
    }
    & \\ \\

    \theutterance \stepcounter{utterance}  
    & & & \multicolumn{4}{p{0.6\linewidth}}{
        \cellcolor[rgb]{0.9,0.9,0.9}{
            \makecell[{{p{\linewidth}}}]{
                \texttt{\tiny{[GM$\langle$P2]}}
                \texttt{BOG: L, (2,2)} \\
            }
        }
    }
     \\ \\

    \theutterance \stepcounter{utterance}  
    & & & \multicolumn{2}{p{0.3\linewidth}}{
        \cellcolor[rgb]{0.9,0.9,0.9}{
            \makecell[{{p{\linewidth}}}]{
                \texttt{\tiny{[GM$|$GM]}}
                \texttt{BOG: L, (2,2)} \\
            }
        }
    }
    & & \\ \\

    \theutterance \stepcounter{utterance}  
    & & & \multicolumn{2}{p{0.3\linewidth}}{
        \cellcolor[rgb]{0.9,0.9,0.9}{
            \makecell[{{p{\linewidth}}}]{
                \texttt{\tiny{[GM$|$GM]}}
                \texttt{Pionós: níl (2,2) folamh, ach tá '├' ann.} \\
\\ 
\texttt{Tá 7 as 8 pionós carntha agaibh le chéile.} \\
\texttt{Déan iarracht arís, le do thoil!} \\
            }
        }
    }
    & & \\ \\

    \theutterance \stepcounter{utterance}  
    & & \multicolumn{4}{p{0.6\linewidth}}{
        \cellcolor[rgb]{0.9,0.9,0.9}{
            \makecell[{{p{\linewidth}}}]{
                \texttt{\tiny{[GM$\rangle$P2]}}
                \texttt{Pionós: níl (2,2) folamh, ach tá '├' ann.} \\
\\ 
\texttt{Tá 7 as 8 pionós carntha agaibh le chéile.} \\
\texttt{Déan iarracht arís, le do thoil!} \\
            }
        }
    }
    & \\ \\

    \theutterance \stepcounter{utterance}  
    & & & \multicolumn{4}{p{0.6\linewidth}}{
        \cellcolor[rgb]{0.9,0.9,0.9}{
            \makecell[{{p{\linewidth}}}]{
                \texttt{\tiny{[GM$\langle$P2]}}
                \texttt{BOG: C, (5,4)} \\
            }
        }
    }
     \\ \\

    \theutterance \stepcounter{utterance}  
    & & & \multicolumn{2}{p{0.3\linewidth}}{
        \cellcolor[rgb]{0.9,0.9,0.9}{
            \makecell[{{p{\linewidth}}}]{
                \texttt{\tiny{[GM$|$GM]}}
                \texttt{BOG: C, (5,4)} \\
            }
        }
    }
    & & \\ \\

    \theutterance \stepcounter{utterance}  
    & & & \multicolumn{2}{p{0.3\linewidth}}{
        \cellcolor[rgb]{0.9,0.9,0.9}{
            \makecell[{{p{\linewidth}}}]{
                \texttt{\tiny{[GM$|$GM]}}
                \texttt{C moved to (5, 4)} \\
            }
        }
    }
    & & \\ \\

    \theutterance \stepcounter{utterance}  
    & & \multicolumn{4}{p{0.6\linewidth}}{
        \cellcolor[rgb]{0.9,0.9,0.9}{
            \makecell[{{p{\linewidth}}}]{
                \texttt{\tiny{[P1$\langle$GM]}}
                \texttt{Cuireadh do theachtaireacht ar aghaidh chuig an imreoir eile.} \\
\\ 
\\ 
\texttt{Tá sibh ag imirt babhta 6 faoi láthair as uasmhéid 12.} \\
\texttt{Tá 7 as 8 pionós carntha agaibh le chéile.} \\
\texttt{Bhog an t{-}imreoir eile réad ar a ngreille.} \\
\\ 
\\ 
\texttt{Cad é do chéad ordú eile?} \\
            }
        }
    }
    & \\ \\

    \theutterance \stepcounter{utterance}  
    & \multicolumn{4}{p{0.6\linewidth}}{
        \cellcolor[rgb]{0.9,0.9,0.9}{
            \makecell[{{p{\linewidth}}}]{
                \texttt{\tiny{[P1$\rangle$GM]}}
                \texttt{ABAIR: An bhfuil tú sásta leis an mbogadh sin? Má bhogadh mé mo 'L' go (2,2), an bhfuil spota folamh agat ag (1,1)?} \\
            }
        }
    }
    & & \\ \\

    \theutterance \stepcounter{utterance}  
    & & & \multicolumn{2}{p{0.3\linewidth}}{
        \cellcolor[rgb]{0.9,0.9,0.9}{
            \makecell[{{p{\linewidth}}}]{
                \texttt{\tiny{[GM$|$GM]}}
                \texttt{ABAIR: An bhfuil tú sásta leis an mbogadh sin? Má bhogadh mé mo 'L' go (2,2), an bhfuil spota folamh agat ag (1,1)?} \\
            }
        }
    }
    & & \\ \\

    \theutterance \stepcounter{utterance}  
    & & \multicolumn{4}{p{0.6\linewidth}}{
        \cellcolor[rgb]{0.9,0.9,0.9}{
            \makecell[{{p{\linewidth}}}]{
                \texttt{\tiny{[GM$\rangle$P2]}}
                \texttt{Bogadh 'C' go (5,4) go rathúil. Seo mar atá do ghreille nuashonraithe:} \\
\texttt{```} \\
\texttt{ 1234567} \\
\texttt{╔═╤═╤═══╗} \\
\texttt{║◌│◌│◌◌◌║ 1} \\
\texttt{║◌├─┤◌L◌║ 2} \\
\texttt{╟─┤P├─┬─╢ 3} \\
\texttt{║◌│◌│C│◌║ 4} \\
\texttt{╟─┼─┘◌│◌║ 5} \\
\texttt{║◌│◌◌┌┴─╢ 6} \\
\texttt{║◌│◌◌│◌◌║ 7} \\
\texttt{╚═╧══╧══╝} \\
\texttt{```} \\
\texttt{Tá sibh ag imirt babhta 7 faoi láthair as uasmhéid 12.} \\
\texttt{Tá 7 as 8 pionós carntha agaibh le chéile.} \\
\texttt{Sheol an t{-}imreoir eile an teachtaireacht seo:} \\
\\ 
\\ 
\texttt{"An bhfuil tú sásta leis an mbogadh sin? Má bhogadh mé mo 'L' go (2,2), an bhfuil spota folamh agat ag (1,1)?"} \\
\\ 
\\ 
\texttt{Cad é do chéad ordú eile?} \\
            }
        }
    }
    & \\ \\

    \theutterance \stepcounter{utterance}  
    & & & \multicolumn{4}{p{0.6\linewidth}}{
        \cellcolor[rgb]{0.9,0.9,0.9}{
            \makecell[{{p{\linewidth}}}]{
                \texttt{\tiny{[GM$\langle$P2]}}
                \texttt{ABAIR: Níl (1,1) folamh, tá '═' ann. Tá (1,1) lán. Tá mo staid: L ag (2,2), C ag (5,4), agus P ag (4,3).} \\
            }
        }
    }
     \\ \\

    \theutterance \stepcounter{utterance}  
    & & & \multicolumn{2}{p{0.3\linewidth}}{
        \cellcolor[rgb]{0.9,0.9,0.9}{
            \makecell[{{p{\linewidth}}}]{
                \texttt{\tiny{[GM$|$GM]}}
                \texttt{ABAIR: Níl (1,1) folamh, tá '═' ann. Tá (1,1) lán. Tá mo staid: L ag (2,2), C ag (5,4), agus P ag (4,3).} \\
            }
        }
    }
    & & \\ \\

    \theutterance \stepcounter{utterance}  
    & & \multicolumn{4}{p{0.6\linewidth}}{
        \cellcolor[rgb]{0.9,0.9,0.9}{
            \makecell[{{p{\linewidth}}}]{
                \texttt{\tiny{[P1$\langle$GM]}}
                \texttt{Cuireadh do theachtaireacht ar aghaidh chuig an imreoir eile.} \\
\\ 
\\ 
\texttt{Tá sibh ag imirt babhta 7 faoi láthair as uasmhéid 12.} \\
\texttt{Tá 7 as 8 pionós carntha agaibh le chéile.} \\
\texttt{Sheol an t{-}imreoir eile an teachtaireacht seo:} \\
\\ 
\\ 
\texttt{"Níl (1,1) folamh, tá '═' ann. Tá (1,1) lán. Tá mo staid: L ag (2,2), C ag (5,4), agus P ag (4,3)."} \\
\\ 
\\ 
\texttt{Cad é do chéad ordú eile?} \\
            }
        }
    }
    & \\ \\

    \theutterance \stepcounter{utterance}  
    & \multicolumn{4}{p{0.6\linewidth}}{
        \cellcolor[rgb]{0.9,0.9,0.9}{
            \makecell[{{p{\linewidth}}}]{
                \texttt{\tiny{[P1$\rangle$GM]}}
                \texttt{ABAIR: Gabh mo leithscéal faoi (1,1)! Tá sé an{-}casta an greille seo. Má tá do staid: L ag (2,2), C ag (5,4), agus P ag (4,3), agus mo staid: L ag (1,1), C ag (1,2), agus P ag (7,7), an bhfuil muid ag iarraidh na rudaí seo a chur ar líne: L ag (2,2), C ag (5,5), agus P ag (4,3)?} \\
            }
        }
    }
    & & \\ \\

    \theutterance \stepcounter{utterance}  
    & & & \multicolumn{2}{p{0.3\linewidth}}{
        \cellcolor[rgb]{0.9,0.9,0.9}{
            \makecell[{{p{\linewidth}}}]{
                \texttt{\tiny{[GM$|$GM]}}
                \texttt{ABAIR: Gabh mo leithscéal faoi (1,1)! Tá sé an{-}casta an greille seo. Má tá do staid: L ag (2,2), C ag (5,4), agus P ag (4,3), agus mo staid: L ag (1,1), C ag (1,2), agus P ag (7,7), an bhfuil muid ag iarraidh na rudaí seo a chur ar líne: L ag (2,2), C ag (5,5), agus P ag (4,3)?} \\
            }
        }
    }
    & & \\ \\

    \theutterance \stepcounter{utterance}  
    & & \multicolumn{4}{p{0.6\linewidth}}{
        \cellcolor[rgb]{0.9,0.9,0.9}{
            \makecell[{{p{\linewidth}}}]{
                \texttt{\tiny{[GM$\rangle$P2]}}
                \texttt{Cuireadh do theachtaireacht ar aghaidh chuig an imreoir eile.} \\
\\ 
\\ 
\texttt{Tá sibh ag imirt babhta 8 faoi láthair as uasmhéid 12.} \\
\texttt{Tá 7 as 8 pionós carntha agaibh le chéile.} \\
\texttt{Sheol an t{-}imreoir eile an teachtaireacht seo:} \\
\\ 
\\ 
\texttt{"Gabh mo leithscéal faoi (1,1)! Tá sé an{-}casta an greille seo. Má tá do staid: L ag (2,2), C ag (5,4), agus P ag (4,3), agus mo staid: L ag (1,1), C ag (1,2), agus P ag (7,7), an bhfuil muid ag iarraidh na rudaí seo a chur ar líne: L ag (2,2), C ag (5,5), agus P ag (4,3)?"} \\
\\ 
\\ 
\texttt{Cad é do chéad ordú eile?} \\
            }
        }
    }
    & \\ \\

    \theutterance \stepcounter{utterance}  
    & & & \multicolumn{4}{p{0.6\linewidth}}{
        \cellcolor[rgb]{0.9,0.9,0.9}{
            \makecell[{{p{\linewidth}}}]{
                \texttt{\tiny{[GM$\langle$P2]}}
                \texttt{ABAIR: Is sprioc maith é sin! Tá muid an{-}gar do dheireadh. Is é seo mo staid: L ag (2,2), C ag (5,4), agus P ag (4,3).} \\
            }
        }
    }
     \\ \\

    \theutterance \stepcounter{utterance}  
    & & & \multicolumn{2}{p{0.3\linewidth}}{
        \cellcolor[rgb]{0.9,0.9,0.9}{
            \makecell[{{p{\linewidth}}}]{
                \texttt{\tiny{[GM$|$GM]}}
                \texttt{ABAIR: Is sprioc maith é sin! Tá muid an{-}gar do dheireadh. Is é seo mo staid: L ag (2,2), C ag (5,4), agus P ag (4,3).} \\
            }
        }
    }
    & & \\ \\

    \theutterance \stepcounter{utterance}  
    & & \multicolumn{4}{p{0.6\linewidth}}{
        \cellcolor[rgb]{0.9,0.9,0.9}{
            \makecell[{{p{\linewidth}}}]{
                \texttt{\tiny{[P1$\langle$GM]}}
                \texttt{Cuireadh do theachtaireacht ar aghaidh chuig an imreoir eile.} \\
\\ 
\\ 
\texttt{Tá sibh ag imirt babhta 8 faoi láthair as uasmhéid 12.} \\
\texttt{Tá 7 as 8 pionós carntha agaibh le chéile.} \\
\texttt{Sheol an t{-}imreoir eile an teachtaireacht seo:} \\
\\ 
\\ 
\texttt{"Is sprioc maith é sin! Tá muid an{-}gar do dheireadh. Is é seo mo staid: L ag (2,2), C ag (5,4), agus P ag (4,3)."} \\
\\ 
\\ 
\texttt{Cad é do chéad ordú eile?} \\
            }
        }
    }
    & \\ \\

    \theutterance \stepcounter{utterance}  
    & \multicolumn{4}{p{0.6\linewidth}}{
        \cellcolor[rgb]{0.9,0.9,0.9}{
            \makecell[{{p{\linewidth}}}]{
                \texttt{\tiny{[P1$\rangle$GM]}}
                \texttt{BOG: C, (1, 2)} \\
            }
        }
    }
    & & \\ \\

    \theutterance \stepcounter{utterance}  
    & & & \multicolumn{2}{p{0.3\linewidth}}{
        \cellcolor[rgb]{0.9,0.9,0.9}{
            \makecell[{{p{\linewidth}}}]{
                \texttt{\tiny{[GM$|$GM]}}
                \texttt{BOG: C, (1, 2)} \\
            }
        }
    }
    & & \\ \\

    \theutterance \stepcounter{utterance}  
    & & & \multicolumn{2}{p{0.3\linewidth}}{
        \cellcolor[rgb]{0.9,0.9,0.9}{
            \makecell[{{p{\linewidth}}}]{
                \texttt{\tiny{[GM$|$GM]}}
                \texttt{Pionós: níl (1,2) folamh, ach tá 'C' ann.} \\
\\ 
\texttt{Tá 8 as 8 pionós carntha agaibh le chéile.} \\
\texttt{Déan iarracht arís, le do thoil!} \\
            }
        }
    }
    & & \\ \\

    \theutterance \stepcounter{utterance}  
    & & \multicolumn{4}{p{0.6\linewidth}}{
        \cellcolor[rgb]{0.9,0.9,0.9}{
            \makecell[{{p{\linewidth}}}]{
                \texttt{\tiny{[P1$\langle$GM]}}
                \texttt{Pionós: níl (1,2) folamh, ach tá 'C' ann.} \\
\\ 
\texttt{Tá 8 as 8 pionós carntha agaibh le chéile.} \\
\texttt{Déan iarracht arís, le do thoil!} \\
            }
        }
    }
    & \\ \\

    \theutterance \stepcounter{utterance}  
    & \multicolumn{4}{p{0.6\linewidth}}{
        \cellcolor[rgb]{0.9,0.9,0.9}{
            \makecell[{{p{\linewidth}}}]{
                \texttt{\tiny{[P1$\rangle$GM]}}
                \texttt{ABAIR: Ó, tá an 'C' fós ansin! Táimid ag teacht ar an deireadh, ach tá an greille seo an{-}dúrsach. Má tá do staid: L ag (2,2), C ag (5,4), agus P ag (4,3), agus mo staid: L ag (1,1), C ag (1,2), agus P ag (7,7), an bhfuil tú sásta má bhogaim mé mo 'P' ó (7,7) go (4,3)?} \\
            }
        }
    }
    & & \\ \\

    \theutterance \stepcounter{utterance}  
    & & & \multicolumn{2}{p{0.3\linewidth}}{
        \cellcolor[rgb]{0.9,0.9,0.9}{
            \makecell[{{p{\linewidth}}}]{
                \texttt{\tiny{[GM$|$GM]}}
                \texttt{ABAIR: Ó, tá an 'C' fós ansin! Táimid ag teacht ar an deireadh, ach tá an greille seo an{-}dúrsach. Má tá do staid: L ag (2,2), C ag (5,4), agus P ag (4,3), agus mo staid: L ag (1,1), C ag (1,2), agus P ag (7,7), an bhfuil tú sásta má bhogaim mé mo 'P' ó (7,7) go (4,3)?} \\
            }
        }
    }
    & & \\ \\

    \theutterance \stepcounter{utterance}  
    & & \multicolumn{4}{p{0.6\linewidth}}{
        \cellcolor[rgb]{0.9,0.9,0.9}{
            \makecell[{{p{\linewidth}}}]{
                \texttt{\tiny{[GM$\rangle$P2]}}
                \texttt{Cuireadh do theachtaireacht ar aghaidh chuig an imreoir eile.} \\
\\ 
\\ 
\texttt{Tá sibh ag imirt babhta 9 faoi láthair as uasmhéid 12.} \\
\texttt{Tá 8 as 8 pionós carntha agaibh le chéile.} \\
\texttt{Sheol an t{-}imreoir eile an teachtaireacht seo:} \\
\\ 
\\ 
\texttt{"Ó, tá an 'C' fós ansin! Táimid ag teacht ar an deireadh, ach tá an greille seo an{-}dúrsach. Má tá do staid: L ag (2,2), C ag (5,4), agus P ag (4,3), agus mo staid: L ag (1,1), C ag (1,2), agus P ag (7,7), an bhfuil tú sásta má bhogaim mé mo 'P' ó (7,7) go (4,3)?"} \\
\\ 
\\ 
\texttt{Cad é do chéad ordú eile?} \\
            }
        }
    }
    & \\ \\

    \theutterance \stepcounter{utterance}  
    & & & \multicolumn{4}{p{0.6\linewidth}}{
        \cellcolor[rgb]{0.9,0.9,0.9}{
            \makecell[{{p{\linewidth}}}]{
                \texttt{\tiny{[GM$\langle$P2]}}
                \texttt{ABAIR: Tá an sprioc sin an{-}mhaith! Má bhogann tú P go (4,3), beidh muid an{-}gar. Is é seo mo staid: L ag (2,2), C ag (5,4), agus P ag (4,3).} \\
            }
        }
    }
     \\ \\

    \theutterance \stepcounter{utterance}  
    & & & \multicolumn{2}{p{0.3\linewidth}}{
        \cellcolor[rgb]{0.9,0.9,0.9}{
            \makecell[{{p{\linewidth}}}]{
                \texttt{\tiny{[GM$|$GM]}}
                \texttt{ABAIR: Tá an sprioc sin an{-}mhaith! Má bhogann tú P go (4,3), beidh muid an{-}gar. Is é seo mo staid: L ag (2,2), C ag (5,4), agus P ag (4,3).} \\
            }
        }
    }
    & & \\ \\

    \theutterance \stepcounter{utterance}  
    & & \multicolumn{4}{p{0.6\linewidth}}{
        \cellcolor[rgb]{0.9,0.9,0.9}{
            \makecell[{{p{\linewidth}}}]{
                \texttt{\tiny{[P1$\langle$GM]}}
                \texttt{Cuireadh do theachtaireacht ar aghaidh chuig an imreoir eile.} \\
\\ 
\\ 
\texttt{Tá sibh ag imirt babhta 9 faoi láthair as uasmhéid 12.} \\
\texttt{Tá 8 as 8 pionós carntha agaibh le chéile.} \\
\texttt{Sheol an t{-}imreoir eile an teachtaireacht seo:} \\
\\ 
\\ 
\texttt{"Tá an sprioc sin an{-}mhaith! Má bhogann tú P go (4,3), beidh muid an{-}gar. Is é seo mo staid: L ag (2,2), C ag (5,4), agus P ag (4,3)."} \\
\\ 
\\ 
\texttt{Cad é do chéad ordú eile?} \\
            }
        }
    }
    & \\ \\

    \theutterance \stepcounter{utterance}  
    & \multicolumn{4}{p{0.6\linewidth}}{
        \cellcolor[rgb]{0.9,0.9,0.9}{
            \makecell[{{p{\linewidth}}}]{
                \texttt{\tiny{[P1$\rangle$GM]}}
                \texttt{BOG: P, (4, 3)} \\
            }
        }
    }
    & & \\ \\

    \theutterance \stepcounter{utterance}  
    & & & \multicolumn{2}{p{0.3\linewidth}}{
        \cellcolor[rgb]{0.9,0.9,0.9}{
            \makecell[{{p{\linewidth}}}]{
                \texttt{\tiny{[GM$|$GM]}}
                \texttt{BOG: P, (4, 3)} \\
            }
        }
    }
    & & \\ \\

    \theutterance \stepcounter{utterance}  
    & & & \multicolumn{2}{p{0.3\linewidth}}{
        \cellcolor[rgb]{0.9,0.9,0.9}{
            \makecell[{{p{\linewidth}}}]{
                \texttt{\tiny{[GM$|$GM]}}
                \texttt{Pionós: níl (4,3) folamh, ach tá '├' ann.} \\
\\ 
\texttt{Tá 9 as 8 pionós carntha agaibh le chéile.} \\
\texttt{Déan iarracht arís, le do thoil!} \\
            }
        }
    }
    & & \\ \\

    \theutterance \stepcounter{utterance}  
    & & \multicolumn{4}{p{0.6\linewidth}}{
        \cellcolor[rgb]{0.9,0.9,0.9}{
            \makecell[{{p{\linewidth}}}]{
                \texttt{\tiny{[P1$\langle$GM]}}
                \texttt{Pionós: níl (4,3) folamh, ach tá '├' ann.} \\
\\ 
\texttt{Tá 9 as 8 pionós carntha agaibh le chéile.} \\
\texttt{Déan iarracht arís, le do thoil!} \\
            }
        }
    }
    & \\ \\

    \theutterance \stepcounter{utterance}  
    & \multicolumn{4}{p{0.6\linewidth}}{
        \cellcolor[rgb]{0.9,0.9,0.9}{
            \makecell[{{p{\linewidth}}}]{
                \texttt{\tiny{[P1$\rangle$GM]}}
                \texttt{ABAIR: Ó, tá an spota (4,3) lán le fráma! Táimid ag teacht ar an 10ú babhta agus tá 9 as 8 pionóis againn. Tá an cluiche an{-}deacair mar tá na líneanna ag claochlú na comhordanáidí. An féidir leat deimhniú go bhfuil spota folamh agat ag (4,4) nó (5,3)?} \\
            }
        }
    }
    & & \\ \\

    \theutterance \stepcounter{utterance}  
    & & & \multicolumn{2}{p{0.3\linewidth}}{
        \cellcolor[rgb]{0.9,0.9,0.9}{
            \makecell[{{p{\linewidth}}}]{
                \texttt{\tiny{[GM$|$GM]}}
                \texttt{ABAIR: Ó, tá an spota (4,3) lán le fráma! Táimid ag teacht ar an 10ú babhta agus tá 9 as 8 pionóis againn. Tá an cluiche an{-}deacair mar tá na líneanna ag claochlú na comhordanáidí. An féidir leat deimhniú go bhfuil spota folamh agat ag (4,4) nó (5,3)?} \\
            }
        }
    }
    & & \\ \\

    \theutterance \stepcounter{utterance}  
    & & \multicolumn{4}{p{0.6\linewidth}}{
        \cellcolor[rgb]{0.9,0.9,0.9}{
            \makecell[{{p{\linewidth}}}]{
                \texttt{\tiny{[GM$\rangle$P2]}}
                \texttt{Cuireadh do theachtaireacht ar aghaidh chuig an imreoir eile.} \\
\\ 
\\ 
\texttt{Tá sibh ag imirt babhta 10 faoi láthair as uasmhéid 12.} \\
\texttt{Tá 9 as 8 pionós carntha agaibh le chéile.} \\
\texttt{Sheol an t{-}imreoir eile an teachtaireacht seo:} \\
\\ 
\\ 
\texttt{"Ó, tá an spota (4,3) lán le fráma! Táimid ag teacht ar an 10ú babhta agus tá 9 as 8 pionóis againn. Tá an cluiche an{-}deacair mar tá na líneanna ag claochlú na comhordanáidí. An féidir leat deimhniú go bhfuil spota folamh agat ag (4,4) nó (5,3)?"} \\
\\ 
\\ 
\texttt{Cad é do chéad ordú eile?} \\
            }
        }
    }
    & \\ \\

    \theutterance \stepcounter{utterance}  
    & & & \multicolumn{4}{p{0.6\linewidth}}{
        \cellcolor[rgb]{0.9,0.9,0.9}{
            \makecell[{{p{\linewidth}}}]{
                \texttt{\tiny{[GM$\langle$P2]}}
                \texttt{ABAIR: Tá (4,4) folamh agus (5,3) folamh freisin. Tá mo staid: L ag (2,2), C ag (5,4), agus P ag (4,3).} \\
            }
        }
    }
     \\ \\

    \theutterance \stepcounter{utterance}  
    & & & \multicolumn{2}{p{0.3\linewidth}}{
        \cellcolor[rgb]{0.9,0.9,0.9}{
            \makecell[{{p{\linewidth}}}]{
                \texttt{\tiny{[GM$|$GM]}}
                \texttt{ABAIR: Tá (4,4) folamh agus (5,3) folamh freisin. Tá mo staid: L ag (2,2), C ag (5,4), agus P ag (4,3).} \\
            }
        }
    }
    & & \\ \\

    \theutterance \stepcounter{utterance}  
    & & \multicolumn{4}{p{0.6\linewidth}}{
        \cellcolor[rgb]{0.9,0.9,0.9}{
            \makecell[{{p{\linewidth}}}]{
                \texttt{\tiny{[P1$\langle$GM]}}
                \texttt{Cuireadh do theachtaireacht ar aghaidh chuig an imreoir eile.} \\
\\ 
\\ 
\texttt{Tá sibh ag imirt babhta 10 faoi láthair as uasmhéid 12.} \\
\texttt{Tá 9 as 8 pionós carntha agaibh le chéile.} \\
\texttt{Sheol an t{-}imreoir eile an teachtaireacht seo:} \\
\\ 
\\ 
\texttt{"Tá (4,4) folamh agus (5,3) folamh freisin. Tá mo staid: L ag (2,2), C ag (5,4), agus P ag (4,3)."} \\
\\ 
\\ 
\texttt{Cad é do chéad ordú eile?} \\
            }
        }
    }
    & \\ \\

    \theutterance \stepcounter{utterance}  
    & \multicolumn{4}{p{0.6\linewidth}}{
        \cellcolor[rgb]{0.9,0.9,0.9}{
            \makecell[{{p{\linewidth}}}]{
                \texttt{\tiny{[P1$\rangle$GM]}}
                \texttt{BOG: P, (4, 4)} \\
            }
        }
    }
    & & \\ \\

    \theutterance \stepcounter{utterance}  
    & & & \multicolumn{2}{p{0.3\linewidth}}{
        \cellcolor[rgb]{0.9,0.9,0.9}{
            \makecell[{{p{\linewidth}}}]{
                \texttt{\tiny{[GM$|$GM]}}
                \texttt{BOG: P, (4, 4)} \\
            }
        }
    }
    & & \\ \\

    \theutterance \stepcounter{utterance}  
    & & & \multicolumn{2}{p{0.3\linewidth}}{
        \cellcolor[rgb]{0.9,0.9,0.9}{
            \makecell[{{p{\linewidth}}}]{
                \texttt{\tiny{[GM$|$GM]}}
                \texttt{Pionós: níl (4,4) folamh, ach tá '│' ann.} \\
\\ 
\texttt{Tá 10 as 8 pionós carntha agaibh le chéile.} \\
\texttt{Déan iarracht arís, le do thoil!} \\
            }
        }
    }
    & & \\ \\

    \theutterance \stepcounter{utterance}  
    & & \multicolumn{4}{p{0.6\linewidth}}{
        \cellcolor[rgb]{0.9,0.9,0.9}{
            \makecell[{{p{\linewidth}}}]{
                \texttt{\tiny{[P1$\langle$GM]}}
                \texttt{Pionós: níl (4,4) folamh, ach tá '│' ann.} \\
\\ 
\texttt{Tá 10 as 8 pionós carntha agaibh le chéile.} \\
\texttt{Déan iarracht arís, le do thoil!} \\
            }
        }
    }
    & \\ \\

    \theutterance \stepcounter{utterance}  
    & \multicolumn{4}{p{0.6\linewidth}}{
        \cellcolor[rgb]{0.9,0.9,0.9}{
            \makecell[{{p{\linewidth}}}]{
                \texttt{\tiny{[P1$\rangle$GM]}}
                \texttt{ABAIR: Tá an cluiche seo an{-}deacair! Táimid ag teacht ar an 11ú babhta agus tá 10 as 8 pionóis againn. Tá na líneanna sa ghreille ag cur isteach ar gach iarracht bogtha. An bhfuil spota folamh agat ag (5,3) nó (4,5)?} \\
            }
        }
    }
    & & \\ \\

    \theutterance \stepcounter{utterance}  
    & & & \multicolumn{2}{p{0.3\linewidth}}{
        \cellcolor[rgb]{0.9,0.9,0.9}{
            \makecell[{{p{\linewidth}}}]{
                \texttt{\tiny{[GM$|$GM]}}
                \texttt{ABAIR: Tá an cluiche seo an{-}deacair! Táimid ag teacht ar an 11ú babhta agus tá 10 as 8 pionóis againn. Tá na líneanna sa ghreille ag cur isteach ar gach iarracht bogtha. An bhfuil spota folamh agat ag (5,3) nó (4,5)?} \\
            }
        }
    }
    & & \\ \\

    \theutterance \stepcounter{utterance}  
    & & \multicolumn{4}{p{0.6\linewidth}}{
        \cellcolor[rgb]{0.9,0.9,0.9}{
            \makecell[{{p{\linewidth}}}]{
                \texttt{\tiny{[GM$\rangle$P2]}}
                \texttt{Cuireadh do theachtaireacht ar aghaidh chuig an imreoir eile.} \\
\\ 
\\ 
\texttt{Tá sibh ag imirt babhta 11 faoi láthair as uasmhéid 12.} \\
\texttt{Tá 10 as 8 pionós carntha agaibh le chéile.} \\
\texttt{Sheol an t{-}imreoir eile an teachtaireacht seo:} \\
\\ 
\\ 
\texttt{"Tá an cluiche seo an{-}deacair! Táimid ag teacht ar an 11ú babhta agus tá 10 as 8 pionóis againn. Tá na líneanna sa ghreille ag cur isteach ar gach iarracht bogtha. An bhfuil spota folamh agat ag (5,3) nó (4,5)?"} \\
\\ 
\\ 
\texttt{Cad é do chéad ordú eile?} \\
            }
        }
    }
    & \\ \\

    \theutterance \stepcounter{utterance}  
    & & & \multicolumn{4}{p{0.6\linewidth}}{
        \cellcolor[rgb]{0.9,0.9,0.9}{
            \makecell[{{p{\linewidth}}}]{
                \texttt{\tiny{[GM$\langle$P2]}}
                \texttt{ABAIR: Tá (5,3) folamh agus (4,5) folamh freisin. Tá mo staid: L ag (2,2), C ag (5,4), agus P ag (4,3).} \\
            }
        }
    }
     \\ \\

    \theutterance \stepcounter{utterance}  
    & & & \multicolumn{2}{p{0.3\linewidth}}{
        \cellcolor[rgb]{0.9,0.9,0.9}{
            \makecell[{{p{\linewidth}}}]{
                \texttt{\tiny{[GM$|$GM]}}
                \texttt{ABAIR: Tá (5,3) folamh agus (4,5) folamh freisin. Tá mo staid: L ag (2,2), C ag (5,4), agus P ag (4,3).} \\
            }
        }
    }
    & & \\ \\

    \theutterance \stepcounter{utterance}  
    & & \multicolumn{4}{p{0.6\linewidth}}{
        \cellcolor[rgb]{0.9,0.9,0.9}{
            \makecell[{{p{\linewidth}}}]{
                \texttt{\tiny{[P1$\langle$GM]}}
                \texttt{Cuireadh do theachtaireacht ar aghaidh chuig an imreoir eile.} \\
\\ 
\\ 
\texttt{Tá sibh ag imirt babhta 11 faoi láthair as uasmhéid 12.} \\
\texttt{Tá 10 as 8 pionós carntha agaibh le chéile.} \\
\texttt{Sheol an t{-}imreoir eile an teachtaireacht seo:} \\
\\ 
\\ 
\texttt{"Tá (5,3) folamh agus (4,5) folamh freisin. Tá mo staid: L ag (2,2), C ag (5,4), agus P ag (4,3)."} \\
\\ 
\\ 
\texttt{Cad é do chéad ordú eile?} \\
            }
        }
    }
    & \\ \\

    \theutterance \stepcounter{utterance}  
    & \multicolumn{4}{p{0.6\linewidth}}{
        \cellcolor[rgb]{0.9,0.9,0.9}{
            \makecell[{{p{\linewidth}}}]{
                \texttt{\tiny{[P1$\rangle$GM]}}
                \texttt{BOG: P, (5, 3)} \\
            }
        }
    }
    & & \\ \\

    \theutterance \stepcounter{utterance}  
    & & & \multicolumn{2}{p{0.3\linewidth}}{
        \cellcolor[rgb]{0.9,0.9,0.9}{
            \makecell[{{p{\linewidth}}}]{
                \texttt{\tiny{[GM$|$GM]}}
                \texttt{BOG: P, (5, 3)} \\
            }
        }
    }
    & & \\ \\

    \theutterance \stepcounter{utterance}  
    & & & \multicolumn{2}{p{0.3\linewidth}}{
        \cellcolor[rgb]{0.9,0.9,0.9}{
            \makecell[{{p{\linewidth}}}]{
                \texttt{\tiny{[GM$|$GM]}}
                \texttt{Pionós: níl (5,3) folamh, ach tá '─' ann.} \\
\\ 
\texttt{Tá 11 as 8 pionós carntha agaibh le chéile.} \\
\texttt{Déan iarracht arís, le do thoil!} \\
            }
        }
    }
    & & \\ \\

    \theutterance \stepcounter{utterance}  
    & & \multicolumn{4}{p{0.6\linewidth}}{
        \cellcolor[rgb]{0.9,0.9,0.9}{
            \makecell[{{p{\linewidth}}}]{
                \texttt{\tiny{[P1$\langle$GM]}}
                \texttt{Pionós: níl (5,3) folamh, ach tá '─' ann.} \\
\\ 
\texttt{Tá 11 as 8 pionós carntha agaibh le chéile.} \\
\texttt{Déan iarracht arís, le do thoil!} \\
            }
        }
    }
    & \\ \\

    \theutterance \stepcounter{utterance}  
    & \multicolumn{4}{p{0.6\linewidth}}{
        \cellcolor[rgb]{0.9,0.9,0.9}{
            \makecell[{{p{\linewidth}}}]{
                \texttt{\tiny{[P1$\rangle$GM]}}
                \texttt{ABAIR: Táimid ag an 12ú babhta agus tá 11 as 8 pionóis againn. Tá an cluiche seo an{-}deacair mar tá na líneanna sa ghreille ag claochlú na spotaí folamacha! Má tá do staid: L ag (2,2), C ag (5,4), agus P ag (4,3), agus mo staid: L ag (1,1), C ag (1,2), agus P ag (5,3), an bhfuil tú sásta má bhogaim mé mo 'P' go (5,4) chun teacht in éineacht leat?} \\
            }
        }
    }
    & & \\ \\

    \theutterance \stepcounter{utterance}  
    & & & \multicolumn{2}{p{0.3\linewidth}}{
        \cellcolor[rgb]{0.9,0.9,0.9}{
            \makecell[{{p{\linewidth}}}]{
                \texttt{\tiny{[GM$|$GM]}}
                \texttt{ABAIR: Táimid ag an 12ú babhta agus tá 11 as 8 pionóis againn. Tá an cluiche seo an{-}deacair mar tá na líneanna sa ghreille ag claochlú na spotaí folamacha! Má tá do staid: L ag (2,2), C ag (5,4), agus P ag (4,3), agus mo staid: L ag (1,1), C ag (1,2), agus P ag (5,3), an bhfuil tú sásta má bhogaim mé mo 'P' go (5,4) chun teacht in éineacht leat?} \\
            }
        }
    }
    & & \\ \\

    \theutterance \stepcounter{utterance}  
    & & \multicolumn{4}{p{0.6\linewidth}}{
        \cellcolor[rgb]{0.9,0.9,0.9}{
            \makecell[{{p{\linewidth}}}]{
                \texttt{\tiny{[GM$\rangle$P2]}}
                \texttt{Cuireadh do theachtaireacht ar aghaidh chuig an imreoir eile.} \\
\\ 
\\ 
\texttt{Tá sibh ag imirt babhta 12 faoi láthair as uasmhéid 12.} \\
\texttt{Tá 11 as 8 pionós carntha agaibh le chéile.} \\
\texttt{Sheol an t{-}imreoir eile an teachtaireacht seo:} \\
\\ 
\\ 
\texttt{"Táimid ag an 12ú babhta agus tá 11 as 8 pionóis againn. Tá an cluiche seo an{-}deacair mar tá na líneanna sa ghreille ag claochlú na spotaí folamacha! Má tá do staid: L ag (2,2), C ag (5,4), agus P ag (4,3), agus mo staid: L ag (1,1), C ag (1,2), agus P ag (5,3), an bhfuil tú sásta má bhogaim mé mo 'P' go (5,4) chun teacht in éineacht leat?"} \\
\\ 
\\ 
\texttt{Cad é do chéad ordú eile?} \\
            }
        }
    }
    & \\ \\

    \theutterance \stepcounter{utterance}  
    & & & \multicolumn{4}{p{0.6\linewidth}}{
        \cellcolor[rgb]{0.9,0.9,0.9}{
            \makecell[{{p{\linewidth}}}]{
                \texttt{\tiny{[GM$\langle$P2]}}
                \texttt{ABAIR: críochnaithe?} \\
            }
        }
    }
     \\ \\

    \theutterance \stepcounter{utterance}  
    & & & \multicolumn{2}{p{0.3\linewidth}}{
        \cellcolor[rgb]{0.9,0.9,0.9}{
            \makecell[{{p{\linewidth}}}]{
                \texttt{\tiny{[GM$|$GM]}}
                \texttt{ABAIR: críochnaithe?} \\
            }
        }
    }
    & & \\ \\

    \theutterance \stepcounter{utterance}  
    & & & \multicolumn{2}{p{0.3\linewidth}}{
        \cellcolor[rgb]{0.9,0.9,0.9}{
            \makecell[{{p{\linewidth}}}]{
                \texttt{\tiny{[GM$|$GM]}}
                \texttt{Maximum number of rounds reached} \\
            }
        }
    }
    & & \\ \\

    \theutterance \stepcounter{utterance}  
    & & & \multicolumn{2}{p{0.3\linewidth}}{
        \cellcolor[rgb]{0.9,0.9,0.9}{
            \makecell[{{p{\linewidth}}}]{
                \texttt{\tiny{[GM$|$GM]}}
                \texttt{Initial states:} \\
\texttt{```} \\
\texttt{Player 1:} \\
\texttt{ 1234567} \\
\texttt{╔═╤═╤═══╗} \\
\texttt{║L│◌│◌◌◌║ 1} \\
\texttt{║C├─┤◌◌◌║ 2} \\
\texttt{╟─┤◌├─┬─╢ 3} \\
\texttt{║◌│◌│◌│◌║ 4} \\
\texttt{╟─┼─┘◌│◌║ 5} \\
\texttt{║◌│◌◌┌┴─╢ 6} \\
\texttt{║◌│◌◌│◌P║ 7} \\
\texttt{╚═╧══╧══╝} \\
\\ 
\\ 
\texttt{Player 2:} \\
\texttt{ 1234567} \\
\texttt{╔═╤═╤═══╗} \\
\texttt{║◌│◌│◌◌◌║ 1} \\
\texttt{║◌├─┤◌L◌║ 2} \\
\texttt{╟─┤P├─┬─╢ 3} \\
\texttt{║◌│◌│◌│◌║ 4} \\
\texttt{╟─┼─┘C│◌║ 5} \\
\texttt{║◌│◌◌┌┴─╢ 6} \\
\texttt{║◌│◌◌│◌◌║ 7} \\
\texttt{╚═╧══╧══╝} \\
\\ 
\texttt{```} \\
            }
        }
    }
    & & \\ \\

    \theutterance \stepcounter{utterance}  
    & & & \multicolumn{2}{p{0.3\linewidth}}{
        \cellcolor[rgb]{0.9,0.9,0.9}{
            \makecell[{{p{\linewidth}}}]{
                \texttt{\tiny{[GM$|$GM]}}
                \texttt{End states:} \\
\texttt{```} \\
\texttt{Player 1:} \\
\texttt{ 1234567} \\
\texttt{╔═╤═╤═══╗} \\
\texttt{║L│◌│◌◌◌║ 1} \\
\texttt{║C├─┤◌◌◌║ 2} \\
\texttt{╟─┤◌├─┬─╢ 3} \\
\texttt{║◌│◌│◌│◌║ 4} \\
\texttt{╟─┼─┘◌│◌║ 5} \\
\texttt{║◌│◌◌┌┴─╢ 6} \\
\texttt{║◌│◌◌│◌P║ 7} \\
\texttt{╚═╧══╧══╝} \\
\\ 
\\ 
\texttt{Player 2:} \\
\texttt{ 1234567} \\
\texttt{╔═╤═╤═══╗} \\
\texttt{║◌│◌│◌◌◌║ 1} \\
\texttt{║◌├─┤◌L◌║ 2} \\
\texttt{╟─┤P├─┬─╢ 3} \\
\texttt{║◌│◌│C│◌║ 4} \\
\texttt{╟─┼─┘◌│◌║ 5} \\
\texttt{║◌│◌◌┌┴─╢ 6} \\
\texttt{║◌│◌◌│◌◌║ 7} \\
\texttt{╚═╧══╧══╝} \\
\\ 
\texttt{```} \\
            }
        }
    }
    & & \\ \\

    \theutterance \stepcounter{utterance}  
    & & & \multicolumn{2}{p{0.3\linewidth}}{
        \cellcolor[rgb]{0.9,0.9,0.9}{
            \makecell[{{p{\linewidth}}}]{
                \texttt{\tiny{[GM$|$GM]}}
                \texttt{* success: True} \\
\texttt{* lose: True} \\
\texttt{* aborted: False} \\
\texttt{{-}{-}{-}{-}{-}{-}{-}} \\
\texttt{* Shifts: 0.00} \\
\texttt{* Max Shifts: 4.00} \\
\texttt{* Min Shifts: 2.00} \\
\texttt{* End Distance Sum: 15.23} \\
\texttt{* Init Distance Sum: 15.76} \\
\texttt{* Expected Distance Sum: 12.57} \\
\texttt{* Penalties: 11.00} \\
\texttt{* Max Penalties: 8.00} \\
\texttt{* Rounds: 12.00} \\
\texttt{* Max Rounds: 12.00} \\
\texttt{* Object Count: 3.00} \\
            }
        }
    }
    & & \\ \\

\end{supertabular}
}

\end{document}
